% sections/00-introduccion.tex
\section*{Introducción}

El presente taller se orienta a la determinación de las actividades involucradas en la explotación de vulnerabilidades sobre aplicaciones web deliberadamente vulnerables, esto mediante la ejecución práctica de ataques de \eng{SQL Injection} (inyección de sentencias SQL) y de \eng{Cross-Site Scripting} (XSS), de forma que se pueda evidenciar, en un entorno controlado, el impacto real que tiene la ausencia de controles de desarrollo seguro sobre aplicaciones expuestas en internet.

En términos generales, las tareas del laboratorio se agrupan en dos frentes:

\begin{itemize}
  \item \term{SQL Injection:} ejecución de ataques de inyección SQL tanto de forma activa (apoyados de la herramienta \texttt{sqlmap} sobre un objetivo identificado mediante \eng{dorks} de búsqueda), como de forma pasiva (sobrepasando un esquema de autenticación con sentencias SQL inyectadas manualmente sobre los campos del formulario de inicio de sesión).
  \item \term{XSS:} resolución de los seis retos del \eng{XSS Game} de Google \cite{xss-game}, incluyendo el nombre de uno de los integrantes del equipo dentro del script inyectado, tal como lo exige la evidencia de la práctica.
\end{itemize}

Como requisito previo se dispuso de una máquina virtual con la distribución Kali Linux (descargable desde \url{https://www.kali.org/downloads/}), dentro de la cual, en la ruta de menú \texttt{APLICACIONES/04} (evaluación de bases de datos), residen las utilidades de auditoría web que se usan a lo largo de la práctica, además de contar de base con las herramientas de línea de comando y navegadores necesarios para el desarrollo de los ejercicios.
